%-----------------------------------
\section{Classificação da pesquisa científica}
\label{sec_classificacao}

Toda pesquisa pode ser classificada segundo diferentes critérios. Os mais usados em trabalhos acadêmicos são a natureza, a abordagem do problema, os objetivos e os procedimentos técnicos \cite{gil2017projetos,marconi2017fundamentos}, sintetizados no Quadro~\ref{qua_classificacao}. Essa classificação deve aparecer no início do capítulo de metodologia (Seção~\ref{sec_met_classificacao}).

\begin{quadro}[htb]
	\centering
	\caption{Critérios de classificação da pesquisa.}
	\label{qua_classificacao}
	\small
	\begin{tabularx}{\textwidth}{|l|l|X|}
		\hline
		\textbf{Critério} & \textbf{Classificação} & \textbf{Descrição} \\ \hline
		\multirow{2}{*}{Natureza}
			& Básica & Gera conhecimento novo, sem aplicação prática prevista. \\ \cline{2-3}
			& Aplicada & Gera conhecimento para resolver um problema específico. \\ \hline
		\multirow{3}{*}{Abordagem}
			& Quantitativa & Usa dados numéricos e análise estatística. \\ \cline{2-3}
			& Qualitativa & Interpreta fenômenos a partir de dados não numéricos. \\ \cline{2-3}
			& Mista & Combina as duas abordagens. \\ \hline
		\multirow{3}{*}{Objetivos}
			& Exploratória & Proporciona maior familiaridade com o problema. \\ \cline{2-3}
			& Descritiva & Descreve características de uma população ou fenômeno. \\ \cline{2-3}
			& Explicativa & Identifica fatores que causam ou contribuem para um fenômeno. \\ \hline
		\multirow{6}{*}{Procedimentos}
			& Bibliográfica & Baseada em material já publicado. \\ \cline{2-3}
			& Documental & Baseada em documentos sem tratamento analítico. \\ \cline{2-3}
			& Experimental & Manipula variáveis e observa seus efeitos. \\ \cline{2-3}
			& Levantamento & Interroga diretamente as pessoas (\textit{survey}). \\ \cline{2-3}
			& Estudo de caso & Estuda em profundidade um ou poucos casos. \\ \cline{2-3}
			& Pesquisa-ação & Pesquisador e participantes atuam juntos na solução do problema. \\ \hline
	\end{tabularx}
	\fonte{Adaptado de \citeonline{gil2017projetos}.}
\end{quadro}

Os tipos de trabalho descritos nas próximas seções combinam esses critérios de formas típicas. Um TCC de desenvolvimento de software, por exemplo, costuma ser uma pesquisa aplicada, de abordagem mista e objetivo exploratório, enquanto uma revisão sistemática é uma pesquisa bibliográfica, geralmente descritiva.
